%! The Quadratic Formula & the Discriminant — Unit 3, Lesson 3.5 — Exit Ticket (Answer Key)
\documentclass[10pt]{article}
\usepackage{algebra2-article}
\usepackage{algebra2-key}   % pulls in -boxes; do NOT also load -boxes
\newcommand{\TallMath}[1]{$\displaystyle #1\rule[-1.4em]{0pt}{3.2em}$}

\begin{document}

\pageheader{Unit 3, Lesson 3.5}{Exit Ticket --- Answer Key}
\namedateperiod

\vspace{0.12in}

No notes. Show your steps; write every answer in simplest form (use $a+bi$ when the roots are complex).

\begin{enumerate}[leftmargin=1.8em, itemsep=14pt]
  \item \textbf{Solve with the quadratic formula.} Identify $a,b,c$, then substitute.
        \[ x^2-4x+13=0 \qquad x=\ans{$2\pm3i$} \]
        \ansline{$a=1,b=-4,c=13$; $b^2-4ac=16-52=-36$; $x=\dfrac{4\pm\sqrt{-36}}{2}=\dfrac{4\pm6i}{2}=2\pm3i$.}

  \item \textbf{Discriminant only --- no solving.} How many roots does $3x^2-5x+1=0$ have, and of what type?
        \[ b^2-4ac=\ans{$13$} \qquad\Rightarrow\qquad \ans{two distinct real (irrational) roots} \]

  \item \textbf{SOL-style multiple choice.} A quadratic equation has discriminant $b^2-4ac=-8$. Which best
        describes its solutions?
        \begin{enumerate}[label=\Alph*., leftmargin=1.9em, itemsep=1pt]
          \item two distinct real solutions
          \item one repeated real solution
          \item two complex-conjugate solutions \textcolor{keyred}{\textbf{$\leftarrow$ correct}}
          \item no solutions of any kind
        \end{enumerate}
        \begin{itemize}[leftmargin=1.4em, nosep]
          \item[] $b^2-4ac<0$ gives two complex-conjugate roots. \textbf{D} is the trap --- a negative discriminant does \emph{not} mean ``no solutions,'' it means \emph{complex} solutions (Lesson 3.4).
        \end{itemize}
\end{enumerate}

\begin{teachernote}
Sort into ``got it / sign slip in $b^2-4ac$ / dropped the $\pm$ or the $i$ / didn't simplify.'' Telltales:
item~1 stopping at $\dfrac{4\pm\sqrt{-36}}{2}$ without converting to $6i$, or writing $\dfrac{4\pm6i}{2}$
without reducing to $2\pm3i$; item~2 a sign error making $b^2-4ac$ come out $\le0$; item~3 choosing
\textbf{D} (the ``no real solution'' habit from 3.3 that 3.4 replaced). If item~3 or the $i$-conversion is
widely missed, reopen Lesson~3.6 with a 60-second discriminant-and-$i$ recap before starting systems.
\end{teachernote}

\end{document}
